\documentclass[aspectratio=169]{beamer}

\mode<presentation> {
  \usetheme{Madrid}
  \usecolortheme{dove}
  %\setbeamertemplate{footline}
  %\setbeamertemplate{footline}[page number]
  %\setbeamertemplate{navigation symbols}{}
}

% ── graphics ───────────────────────────────────────────────────────────────────
\usepackage{graphicx}
\usepackage{graphics}
\usepackage{wrapfig}
\usepackage{float}
\usepackage{sidecap}
\usepackage{caption}
\usepackage{subcaption}
\graphicspath{{figures/}}

% ── tables & layout ────────────────────────────────────────────────────────────
\usepackage{multicol}
\usepackage{multirow}
\usepackage{array}

% ── math ───────────────────────────────────────────────────────────────────────
\usepackage{amsmath}
\usepackage{mathtools}
%\usepackage{slashed}   % texlive-science
%\usepackage{siunitx}   % texlive-science

% ── links ──────────────────────────────────────────────────────────────────────
\usepackage{hyperref}
\hypersetup{
  colorlinks=true,
  linkcolor=olive,
  filecolor=cyan,
  urlcolor=magenta
}

% ── code listings ──────────────────────────────────────────────────────────────
\usepackage{listings}
\usepackage{xcolor}

\lstset{
    basicstyle=\ttfamily\footnotesize,
    frame=single,
    breaklines=true,
    showstringspaces=false,
}

% ── diagrams & boxes ───────────────────────────────────────────────────────────
\usepackage{tcolorbox}
\usepackage{tikz}
\usetikzlibrary{tikzmark,overlay-beamer-styles,babel}
\usepackage{forest}
\usepackage{eso-pic}
\usepackage[normalem]{ulem}
\usepackage{appendixnumberbeamer}

% ── title info ─────────────────────────────────────────────────────────────────
\title[Git --- Works for You]{Git Works for You}
%% \subtitle{}
\author[Mondal]{Suryanarayan Mondal}
\institute[\href{mailto:suryamondal@gmail.com}{suryamondal@gmail.com}] {
  %% \small{Your Institute}
}
\date{}

% ══════════════════════════════════════════════════════════════════════════════
\begin{document}

\begin{frame}
    \titlepage
\end{frame}

% ─────────────────────────────────────────────────────────────────────────────
\section{The Problem}
% ─────────────────────────────────────────────────────────────────────────────

\begin{frame}[fragile]{Has your folder ever looked like this?}
    \begin{columns}
        \column{0.5\textwidth}
            \begin{lstlisting}
analysis_v1.py
analysis_v2.py
analysis_v2_fixed.py
analysis_final.py
analysis_final_REAL.py
analysis_final_REAL_submitted.py
analysis_final_REAL_submitted_2.py
            \end{lstlisting}
        \column{0.5\textwidth}
            \pause
            \begin{block}{The hard questions}
                \begin{itemize}
                    \item Which one did I submit?
                    \item What changed between v2 and final?
                    \item Which one runs on the cluster?
                    \item What did I change last Tuesday?
                \end{itemize}
            \end{block}
    \end{columns}
    \vspace{1em}
    \pause
    \centering
    \textbf{If this looks familiar --- you were already doing version control.}\\[4pt]
    \hspace{0.5\textwidth}\textbf{Just doing it badly.}
\end{frame}

% ─────────────────────────────────────────────────────────────────────────────

\begin{frame}{You invented version control (without knowing it)}
    \begin{columns}
        \column{0.48\textwidth}
            \begin{block}{What you do manually}
                \begin{tabular}{ll}
                    \texttt{\_v1}, \texttt{\_v2}      & $\leftarrow$ version tag    \\[4pt]
                    \texttt{\_fixed}                   & $\leftarrow$ commit message \\[4pt]
                    \texttt{\_final\_REAL}             & $\leftarrow$ release tag    \\[4pt]
                    keeping old files                  & $\leftarrow$ history        \\
                \end{tabular}
            \end{block}
        \column{0.48\textwidth}
            \pause
            \begin{block}{What git does automatically}
                \begin{tabular}{ll}
                    version number     & assigned           \\[4pt]
                    commit message     & you write it once  \\[4pt]
                    tags               & natively supported \\[4pt]
                    full history       & always, forever    \\
                \end{tabular}
            \end{block}
    \end{columns}
    \vspace{1.2em}
    \pause
    \centering
    \textbf{Git does exactly what you were doing ---\\
    without cluttering your folder.}\\[6pt]
    \pause
    \textit{And with added functionality to make our life much, much easier.}
\end{frame}

% ─────────────────────────────────────────────────────────────────────────────

\begin{frame}{The wheel was already invented}
    \centering
    \vspace{0.5em}
    People have had this problem since the first line of code was written.\\
    \vspace{1em}
    \pause
    Over decades, they collectively built a solution.\\
    It is \textbf{free}. \; It is \textbf{open source}. \; It runs \textbf{everywhere}.\\
    \vspace{1.2em}
    \vspace{0.8em}
    \pause
    Every tool asks you to learn some commands --- that is unavoidable.\\[4pt]
    It becomes muscle memory faster than you think.\\[4pt]
    \pause
    \textbf{Of all the tools you will ever use, git has the fewest commands to learn.}
\end{frame}

% ─────────────────────────────────────────────────────────────────────────────
\section{What is git}
% ─────────────────────────────────────────────────────────────────────────────

\begin{frame}[fragile]{git is just a folder that remembers}
    \begin{columns}
        \column{0.52\textwidth}
            \begin{lstlisting}
my_project/
|-- .git/       <- git lives HERE
|-- analysis.py <- tracked by git
|-- results.csv <- NOT tracked by git
`-- README.md   <- tracked by git
            \end{lstlisting}
            \vspace{0.5em}
            \begin{itemize}
                \item No internet required
                \item No account required
                \item A hidden folder inside your project
                \item Everything git knows lives right there
            \end{itemize}
        \column{0.44\textwidth}
    \end{columns}
    \pause
    \begin{tikzpicture}[remember picture, overlay]
        \node[draw=red, line width=2pt, fill=white, rounded corners=4pt,
              text width=0.55\textwidth, align=center, inner sep=14pt]
              at (current page.center) {
            \textbf{\large That is all git is}\\[8pt]
            A tool installed on your computer.\\[6pt]
            It tracks changes to files you tell it to track.\\[6pt]
            It works completely offline.
        };
    \end{tikzpicture}
\end{frame}

% ─────────────────────────────────────────────────────────────────────────────

\begin{frame}{Two different things. One name.}
    \begin{columns}
        \column{0.48\textwidth}
            \begin{block}{\texttt{git} --- the tool}
                \begin{itemize}
                    \item Runs on your computer
                    \item Works offline
                    \item Free and open source
                    \item No account needed
                    \item History lives in \texttt{.git/}
                \end{itemize}
            \end{block}
        \column{0.48\textwidth}
            \begin{block}{GitHub --- the website}
                \begin{itemize}
                    \item A website in the cloud
                    \item Stores a \textit{copy} of your repo
                    \item Adds a web interface
                    \item Adds collaboration tools
                    \item Optional --- but very useful
                \end{itemize}
            \end{block}
    \end{columns}
    \vspace{1em}
    \pause
    \centering
    \texttt{[Your laptop \;\; .git/]}
    \quad $\overset{\texttt{push}}{\longrightarrow}$ \quad
    \texttt{[GitHub]}\\[6pt]
    \small
    You need \texttt{git} to use GitHub.\\
    You do \textbf{not} need GitHub to use \texttt{git}.
\end{frame}

% ─────────────────────────────────────────────────────────────────────────────
\section{How git works}
% ─────────────────────────────────────────────────────────────────────────────

\begin{frame}{How git works --- four commands}
    \begin{columns}
        \column{0.48\textwidth}
            \begin{block}{\texttt{git init}}
                Tell git to start tracking\\
                this directory.\\[4pt]
                \texttt{\$ git init}\\[4pt]
                \textit{Creates \texttt{.git/}. Done once per project.}
            \end{block}
            \vspace{0.4em}
            \uncover<4->{
            \begin{block}{\texttt{git push}}
                Send your commits to the remote.\\[4pt]
                \texttt{\$ git push}\\[4pt]
                \textit{Now it leaves your machine. Now it is safe.}
            \end{block}
            }
        \column{0.48\textwidth}
            \uncover<2->{
            \begin{block}{\texttt{git add}}
                Tell git which files to include\\
                in the next snapshot.\\[4pt]
                \texttt{\$ git add analysis.py README.md}\\[4pt]
                \textit{Completely offline. Nothing moves.}
            \end{block}
            }
            \vspace{0.4em}
            \uncover<3->{
            \begin{block}{\texttt{git commit}}
                Save that snapshot locally\\
                with a short description.\\[4pt]
                \texttt{\$ git commit -m "added new plot"}\\[4pt]
                \textit{Still offline. Lives in \texttt{.git/} only.}
            \end{block}
            }
    \end{columns}
    \uncover<5->{
    \centering
    \vspace{0.4em}
    \texttt{init} $\rightarrow$ \texttt{add} $\rightarrow$ \texttt{commit} $\rightarrow$ \texttt{push} \quad --- \quad \small That is the entire workflow.
    }
\end{frame}

% ─────────────────────────────────────────────────────────────────────────────
\section{What git gives you}
% ─────────────────────────────────────────────────────────────────────────────

\begin{frame}{What git gives YOU --- not the team}
    \begin{columns}
        \column{0.48\textwidth}
            \begin{block}{\emph{Ownership with a timestamp} $\triangleright$}
                Every commit is signed, dated,\\
                and attributed to you.\\[4pt]
                \textit{Your work diary --- automatic.}
            \end{block}
            \vspace{0.5em}
            \begin{block}{\emph{A personal time machine} $\triangleright$}
                Revert anything, anytime.\\
                Experiment without fear.\\[4pt]
                \textit{Nothing is ever truly deleted.}
            \end{block}
        \column{0.48\textwidth}
            \begin{block}{\emph{Work from anywhere} $\triangleright$}
                Office, home, cluster.\\
                One sync command and you are up to date.\\[4pt]
                \textit{No USB. No rsync. No email to yourself.}
            \end{block}
            \vspace{0.5em}
            \begin{block}{\emph{Freedom to break things} $\triangleright$}
                Create an isolated copy of your work.\\
                Try something crazy. Throw it away if it fails.\\[4pt]
                \textit{Your main code is untouched.}
            \end{block}
    \end{columns}
\end{frame}

% ─────────────────────────────────────────────────────────────────────────────
\section{Why push}
% ─────────────────────────────────────────────────────────────────────────────

\begin{frame}{Why push? \; Not just commit locally?}
    \begin{columns}
        \column{0.5\textwidth}
            \begin{alertblock}{Local commits die with your laptop}
                A commit saved only locally disappears if:
                \begin{itemize}
                    \item Hard drive fails
                    \item Laptop is stolen
                    \item Accidental \texttt{rm -rf}
                \end{itemize}
            \end{alertblock}
        \column{0.5\textwidth}
            \pause
            \begin{block}{A push is different}
                \begin{itemize}
                    \item Your work lives on a server
                    \item Timestamped --- publicly
                    \item \textbf{Nobody can dispute\\who wrote it first}
                    \item One push costs less data\\than sending one email
                \end{itemize}
            \end{block}
    \end{columns}
    \vspace{1em}
    \pause
    \centering
    \textbf{Push once or twice a day.\\
    No bandwidth concern. No data concern.\\
    There is only upside.}
\end{frame}

% ─────────────────────────────────────────────────────────────────────────────

\begin{frame}[fragile]{The daily habit}
    \centering
    \begin{block}{Morning --- get what changed}
        \begin{lstlisting}[language=bash]
git pull
        \end{lstlisting}
    \end{block}
    \vspace{0.6em}
    \begin{block}{Evening --- save your day}
        \begin{lstlisting}[language=bash]
git add -u
git commit -m "what I did today"
git push
        \end{lstlisting}
    \end{block}
    \vspace{0.8em}
    \pause
    \textbf{Three commands. \; Thirty seconds. \; Every day.}\\[6pt]
    \small Your laptop can burn tonight. \; Your work is safe.
\end{frame}

% ─────────────────────────────────────────────────────────────────────────────
\section{Summary}
% ─────────────────────────────────────────────────────────────────────────────

\begin{frame}[fragile]{Summary}
    \begin{enumerate}
        \setlength\itemsep{0.5em}
        \item \textbf{You were already doing it} ---
              \texttt{\_v1}, \texttt{\_v2} is manual version control
        \item \textbf{git does it properly} ---
              automatic, searchable, always recoverable
        \item \textbf{git is local} ---
              no internet, no account, just a folder
        \item \textbf{GitHub is the bonus} ---
              a copy in the cloud, with a nice face
        \item \textbf{Push daily} ---
              backup $+$ ownership $+$ peace of mind
    \end{enumerate}
    \vspace{1em}
    \pause
    \begin{block}{Start today}
        \begin{lstlisting}[language=bash]
git init
git add your_file.py
git commit -m "start tracking my work"
git push
        \end{lstlisting}
    \end{block}
    \centering
    \vspace{0.4em}
    \textit{One repo. One push. Today.}
\end{frame}

% ─────────────────────────────────────────────────────────────────────────────

\begin{frame}{}
    \centering
    \vspace{2em}
    {\Large Thank you}\\
    \vspace{2em}
    {\large Questions?}\\
    \vspace{2em}
    \small
    Slides and git cheatsheet available at:\\[4pt]
    \texttt{github.com/suryamondal/useful\_git\_commands}
\end{frame}

\end{document}
